% F085｜真实数学图；代表函数不构成抽象解的唯一性证明。
\begin{center}
\begin{minipage}{.49\linewidth}\centering
\begin{tikzpicture}
\begin{axis}[width=6.1cm,height=4.8cm,axis lines=middle,
 xmin=-0.2,xmax=3.5,ymin=-0.5,ymax=10.8,
 xtick={0,1,2,3},xticklabels=\empty,ytick={0,1,4,9},xlabel={$t$},ylabel={$f(t)$},
 title={幂：乘法作用于横坐标与高度},title style={font=\small,yshift=3pt},
 tick label style={font=\small},label style={font=\small},
 ylabel style={at={(axis description cs:0,1)},anchor=south east,xshift=-18pt},
 clip=false,grid=major,grid style={black!8},enlargelimits=false,
 every axis plot/.append style={thick},samples=161]
\addplot[CExplanation,very thick,domain=0:3.1] {x^2};
\draw[dashed,black!35] (axis cs:1,0) |- (axis cs:0,1);
\draw[dashed,black!35] (axis cs:2,0) |- (axis cs:0,4);
\draw[dashed,black!35] (axis cs:3,0) |- (axis cs:0,9);
\addplot[only marks,mark=*,mark size=2pt,orange!80!black] coordinates {(1,1) (2,4) (3,9)};

\node[font=\small,anchor=north,yshift=-3pt] at (axis cs:1,-0.5) {$1$};
\node[font=\small,anchor=north,yshift=-3pt] at (axis cs:2,-0.5) {$2$};
\node[font=\small,anchor=north,yshift=-3pt] at (axis cs:3,-0.5) {$3$};

\end{axis}
\end{tikzpicture}
\end{minipage}\hfill
\begin{minipage}{.49\linewidth}\centering
\begin{tikzpicture}
\begin{axis}[width=6.1cm,height=4.8cm,axis lines=middle,
 xmin=-0.2,xmax=3.5,ymin=-0.4,ymax=7.5,
 xtick={0,1,2,3},xticklabels=\empty,ytick={0,2,4,6},xlabel={$t$},ylabel={$f(t)$},
 title={正比例：加法作用于横坐标与高度},title style={font=\small,yshift=3pt},
 tick label style={font=\small},label style={font=\small},
 ylabel style={at={(axis description cs:0,1)},anchor=south east,xshift=-18pt},
 clip=false,grid=major,grid style={black!8},enlargelimits=false,
 every axis plot/.append style={thick},samples=161]
\addplot[CExplanation,very thick,domain=0:3.1] {2*x};
\draw[dashed,black!35] (axis cs:1,0) |- (axis cs:0,2);
\draw[dashed,black!35] (axis cs:2,0) |- (axis cs:0,4);
\draw[dashed,black!35] (axis cs:3,0) |- (axis cs:0,6);
\addplot[only marks,mark=*,mark size=2pt,orange!80!black] coordinates {(1,2) (2,4) (3,6)};

\node[font=\small,anchor=north,yshift=-3pt] at (axis cs:1,-0.4) {$1$};
\node[font=\small,anchor=north,yshift=-3pt] at (axis cs:2,-0.4) {$2$};
\node[font=\small,anchor=north,yshift=-3pt] at (axis cs:3,-0.4) {$3$};

\end{axis}
\end{tikzpicture}
\end{minipage}
\par\small 左图 $f(\sqrt2\cdot\sqrt2)=f(\sqrt2)^2=4$；右图 $f(1+2)=f(1)+f(2)=6$。
\end{center}
